\section{仿真数据质量：sim-to-real gap 怎么处理}

仿真是机器人数据金字塔中最关键的中间层。它能规模化生成场景、控制变量、复现失败、构造评测。但仿真数据只有在与真实世界差距可控时才有价值。否则模型会学会仿真里的捷径，而不是现实中的能力。

\subsection{仿真质量的四个维度}

\noindent\begin{tabularx}{\textwidth}{p{0.22\textwidth}p{0.34\textwidth}X}
\toprule
维度 & 要求 & 失败表现 \\
\midrule
物理一致性 & 接触、摩擦、重力、动力学可信 & 现实中动作失效。 \\
视觉真实度 & 材质、光照、遮挡、相机噪声合理 & 感知模型过拟合仿真画面。 \\
任务多样性 & 覆盖长尾、异常和复杂组合 & 只会标准场景。 \\
评测可迁移 & 仿真评测能预测真实表现 & 仿真榜单和真实部署脱节。 \\
\bottomrule
\end{tabularx}

\begin{warningbox}{仿真不是越真实越好}
过度追求视觉真实可能成本高但收益低。仿真要服务训练和评测目标：哪些变量影响策略，哪些变量只影响外观，必须分清。
\end{warningbox}

\subsection{本章小结}

仿真数据的质量，不在于“看起来像不像”，而在于它是否能提升真实任务表现，并帮助评测真实能力。

